\setcounter{section}{2}


  \zwischenueberschrift{Übungsaufgaben}


\inputaufgabe


{}


{


Betrachte in


\mathl{{ {\mathbb A}_{ K }^{  3   } }}{} die beiden Ebenen


\mathdisp {E_1=V(3x+4y+5z) \text{ und } E_2 =V(2x-y+3z)} { . }


Parametrisiere den Schnitt


\mathl{E_1 \cap E_2}{.}


}


{} {}


\inputaufgabegibtloesung


{}


{


Bestimme die Schnittpunkte des Einheitskreises mit der Geraden, die durch die beiden Punkte
\mathkor {} {(-1,1)} {und} {(4,-2)} {}
verläuft.


}


{} {}


\inputaufgabe


{}


{


Bestimme die Koordinaten der beiden Schnittpunkte der Geraden $G$ und des Kreises $K$, wobei $G$ durch die Gleichung


\mavergleichskette


{\vergleichskette


{ 3y-4x+2
}


{ = }{ 0
}


{  }{
}


{    }{
}


{   }{
}


}
{}{}{}
und $K$ durch den Mittelpunkt $(2,5)$ und den Radius $7$ gegeben sind.


}


{} {}


\inputaufgabe


{}


{


Berechne die Schnittpunkte der beiden Kurven


\mathdisp {C=V \left(  x^2+2y^2+3xy+x-2  \right)  \text{ und } L=V \left(  4x+3y-5  \right)} { . }


}


{} {}


\inputaufgabe


{}


{


Zeige: Der Durchschnitt von zwei verschiedenen Kreisen in der affinen Ebene ist der Durchschnitt eines Kreises mit einer Geraden.


}


{} {}


\inputaufgabegibtloesung


{}


{


Bestimme die Schnittpunkte des Einheitskreises $E$  mit dem Kreis $K$, der den Mittelpunkt $(1,0)$ und den Radius $2$ besitzt.


}


{} {}


\inputaufgabegibtloesung


{}


{


Bestimme die Schnittpunkte der beiden Ellipsen


\mathdisp {\left\{ (x,y) \in\R^2  \mid   x^2+xy+3y^2 = 3         \right\} \text{   und } \left\{ (x,y) \in\R^2  \mid   2x^2-xy+y^2 = 4         \right\}} { . }


}


{} {}


\inputaufgabegibtloesung


{}


{


Bestimme die Schnittpunkte des Einheitskreises mit der Standardparabel.


}


{} {}


\inputaufgabegibtloesung


{}


{


Es sei


\mavergleichskettedisp


{\vergleichskette


{P
}


{ =} { \left\{ (x,y) \in \R^2  \mid   y = x^2         \right\}
}


{ } {
}


{ } {
}


{ } {
}


}
{}{}{}
die Standardparabel und $K$ der Kreis mit dem Mittelpunkt $(0,1)$ und dem Radius $1$.
\aufzaehlungfuenf{Skizziere
\mathkor {} {P} {und} {K} {.}
}{Erstelle eine Gleichung für $K$.
}{Bestimme die Schnittpunkte


\mathdisp {P \cap K} { . }


}{Beschreibe die untere Kreisbogenhälfte als Graph einer Funktion von


\mathl{[-1,1]}{} nach $\R$.
}{Bestimme, wie die Parabel relativ zum unteren Kreisbogen verläuft.
}


}


{} {}


\inputaufgabe


{}


{


Bestimme alle simultanen Lösungen der beiden Gleichungen


\mathdisp {x^3+y^2 = 2 \text{ und } 2xy =3} {  }


für die
\definitionsverweis {Körper}{}{}


\mavergleichskette


{\vergleichskette


{ K
}


{ = }{ \Z/( 3 )
}


{  }{
}


{    }{
}


{   }{
}


}
{}{}{,}
$\Z/( 5 )$ und $\Z/( 7 )$.


}


{} {}


\inputaufgabegibtloesung


{}


{


Wir betrachten die Varietät der kommutierenden
$2 \times 2$-\definitionsverweis {Matrizen}{}{,}
also die Menge der Matrizenpaare


\mavergleichskettedisp


{\vergleichskette


{ V
}


{ =} { \left\{ (A,B)  \mid   A,B \in \operatorname{Mat}_{  2   } (K)  , \,  AB = BA       \right\}
}


{ \subseteq} { \operatorname{Mat}_{  2   } (K) \times \operatorname{Mat}_{  2   } (K)
}


{ \cong} { { {\mathbb A}_{ K }^{  8   } }
}


{ } {
}


}
{}{}{.}
\aufzaehlungdrei{Zeige, dass dies eine
\definitionsverweis {affine Varietät}{}{}
ist, und bestimme möglichst einfache Gleichungen, die diese Varietät beschreiben.
}{Zeige, dass die Abbildung
\maabbeledisp {\pi} { V } { \operatorname{Mat}_{  2   } (K)
} {(A,B)} { A
} {,}
\definitionsverweis {surjektiv}{}{}
ist.
}{Bestimme das Urbild von


\mathl{\begin{pmatrix}  1   &   1  \\  0  &  0 \end{pmatrix}}{} unter $\pi$.
}


}


{} {}


\inputaufgabe


{}


{


Zeige, dass zu einem Punkt


\mavergleichskette


{\vergleichskette


{ P
}


{ = }{ \left(  a_1   , \,   \ldots  , \,   a_n          \right)
}


{ \in }{ { {\mathbb A}_{  K  }^{  n   } }
}


{    }{
}


{   }{
}


}
{}{}{}
das zugehörige Ideal


\mathdisp {\left(  X_1-a_1, X_2-a_2 , \ldots ,  X_n-a_n  \right)} {  }


\definitionsverweis {maximal}{}{}
ist.


}


{} {}


\inputaufgabe


{}


{


Es sei $K$ ein
\definitionsverweis {algebraisch abgeschlossener Körper}{}{}
und


\mavergleichskette


{\vergleichskette


{ n
}


{ \geq }{ 2
}


{  }{
}


{    }{
}


{   }{
}


}
{}{}{.}
Zeige, dass ein Punkt


\mavergleichskette


{\vergleichskette


{ P
}


{ \in }{ { {\mathbb A}_{  K }^{  n   } }
}


{  }{
}


{    }{
}


{   }{
}


}
{}{}{}
nicht die
\definitionsverweis {Nullstellenmenge}{}{}
zu einem einzigen Polynom ist.


}


{} {}


\inputaufgabe


{}


{


Es sei $K$ ein
\definitionsverweis {endlicher Körper}{}{}
und


\mavergleichskette


{\vergleichskette


{ M
}


{ = }{ \{ P_1,P_2 , \ldots , P_m \}
}


{ \subseteq }{ { {\mathbb A}_{  K  }^{  n   } }
}


{    }{
}


{   }{
}


}
{}{}{}
eine endliche Menge von Punkten. Zeige, dass $M$ die
\definitionsverweis {Nullstellenmenge}{}{}
eines einzigen Polynoms ist.


}


{} {}


\inputaufgabegibtloesung


{}


{


Es sei $K$ ein
\definitionsverweis {Körper}{}{.}
Zeige, dass in


\mathl{K[X,Y,Z,W]}{} die drei Ideale


\mavergleichskettedisp


{\vergleichskette


{ {\mathfrak a}
}


{ =} { \left(  Z^2+W^2 -1, \, X - 2Z^2 +1 ,\, Y-2ZW  \right)
}


{ } {
}


{ } {
}


{ } {
}


}
{}{}{,}


\mavergleichskettedisp


{\vergleichskette


{ {\mathfrak b}
}


{ =} { \left(  Z^2+W^2 -1 ,\, X^2+Y^2 -1 ,\, (1-X)Z -YW,\, YZ- (1+X) W  \right)
}


{ } {
}


{ } {
}


{ } {
}


}
{}{}{}
und


\mavergleichskettedisp


{\vergleichskette


{ {\mathfrak c}
}


{ =} { \left(  Z^2+W^2 -1 ,\, (1-X)Z -YW,\, YZ- (1+X) W  \right)
}


{ } {
}


{ } {
}


{ } {
}


}
{}{}{}
übereinstimmen.


}


{} {}


\inputaufgabegibtloesung


{}


{


Wir betrachten das Polynom


\mavergleichskettedisp


{\vergleichskette


{F
}


{ =} { X^4Y^2+X^2Y^4 -3X^2Y^2+1
}


{ } {
}


{ } {
}


{ } {
}


}
{}{}{.}
\aufzaehlungdreiabc{Finde eine reelle Nullstelle von $F$.
}{Bestätige die Gleichung


\mavergleichskettealigndrucklinks


{\vergleichskettealigndrucklinks


{ F \cdot \left(  X^2+Y^2+1  \right)
}


{ =} { \zeilemitteil { \left(  X^2Y-Y  \right)^2 + \left(  XY^2-X  \right)^2 + \left(  X^2Y^2-1  \right)^2 + { \frac{   1  }{  4  } } \left(  XY^3-X^3Y  \right)^2 } {+ { \frac{   3  }{  4  } } \left(  XY^3+X^3Y-2XY  \right)^2} }


{ } { \zeilemitteil {
} {} }


{ } { \zeilemitteil {
} {} }


{ } { \zeilemitteil {
} {} }


}
{}{}{.}
}{Es sei


\mavergleichskette


{\vergleichskette


{ K
}


{ = }{ \R
}


{  }{
}


{    }{
}


{   }{
}


}
{}{}{.}
Folgere, dass


\mavergleichskettedisp


{\vergleichskette


{ F(x,y)
}


{ \geq} { 0
}


{ } {
}


{ } {
}


{ } {
}


}
{}{}{}
ist für alle Punkte


\mavergleichskette


{\vergleichskette


{ (x,y)
}


{ \in }{ {\mathbb A}^{2}_{ \R }
}


{  }{
}


{    }{
}


{   }{
}


}
{}{}{.}
}


}


{Bemerkung: Nach einem Satz von Artin
\zusatzklammer {Lösung des 17. Hilbertschen Problems} {} {}
kann man jedes reelle Polynom, das nirgendwo negative Werte annimmt, als eine Summe von Quadraten von rationalen Funktionen schreiben. Das vorstehende \stichwort {Motzkin-Polynom} {} $F$ gibt ein konkretes Beispiel dafür, dass man ein solches Polynom im Allgemeinen nicht als Summe von Quadraten von Polynomen schreiben kann. Wir haben aber lediglich die Nichtnegativität bewiesen.} {}


\inputaufgabe


{}


{


Bestimme
\definitionsverweis {Idealerzeuger}{}{}
für ein
\definitionsverweis {Ideal}{}{}


\mavergleichskette


{\vergleichskette


{ {\mathfrak a}
}


{ \subseteq }{ \R [X,Y,Z]
}


{  }{
}


{    }{
}


{   }{
}


}
{}{}{,}
dessen Nullstellenmenge genau die vier Punkte


\mavergleichskettedisp


{\vergleichskette


{ (2,3,4), (1,1/5,0), (0,0,1), (-1,-2,\sqrt{3})
}


{ \in} { { {\mathbb A}_{ \R }^{  3   } }
}


{ } {
}


{ } {
}


{ } {
}


}
{}{}{}
sind.


}


{} {}


\inputaufgabe


{}


{


Es seien


\mavergleichskette


{\vergleichskette


{ {\mathfrak a}, {\mathfrak b}
}


{ \subseteq }{ R
}


{  }{
}


{    }{
}


{   }{
}


}
{}{}{}
\definitionsverweis {Ideale}{}{}
in einem
\definitionsverweis {kommutativen Ring}{}{}
$R$. Zeige, dass die Beziehung


\mavergleichskettedisp


{\vergleichskette


{ {\mathfrak a} \cdot {\mathfrak b}
}


{ \subseteq} { {\mathfrak a} \cap {\mathfrak b}
}


{ } {
}


{ } {
}


{ } {
}


}
{}{}{}
gilt.


}


{} {}


\inputaufgabe


{}


{


Zeige, dass das
\definitionsverweis {Produkt}{}{}
von
\definitionsverweis {Hauptidealen}{}{}
wieder ein Hauptideal ist.


}


{} {}


\inputaufgabegibtloesung


{}


{


Es seien
\mathkor {} {I} {und} {J} {}
\definitionsverweis {Ideale}{}{}
in einem
\definitionsverweis {kommutativen Ring}{}{}
$R$ und sei


\mavergleichskette


{\vergleichskette


{ n
}


{ \in }{ \N
}


{  }{
}


{    }{
}


{   }{
}


}
{}{}{.}
Zeige die Gleichheit


\mavergleichskettedisp


{\vergleichskette


{ (I+J)^n
}


{ =} { I^n + I^{n-1}J+ I^{n-2}J^2 + \cdots + I^2J^{n-2} + IJ^{n-1} +J^n
}


{ } {
}


{ } {
}


{ } {
}


}
{}{}{.}


}


{} {}


\inputaufgabe


{}


{


Es sei $R$ ein
\definitionsverweis {kommutativer Ring}{}{}
mit
\definitionsverweis {Idealen}{}{}


\mavergleichskette


{\vergleichskette


{ {\mathfrak a}, {\mathfrak b}
}


{ \subseteq }{ R
}


{  }{
}


{    }{
}


{   }{
}


}
{}{}{.}
Es sei


\mavergleichskette


{\vergleichskette


{ S
}


{ = }{ R/{\mathfrak b}
}


{  }{
}


{    }{
}


{   }{
}


}
{}{}{}
und


\mavergleichskette


{\vergleichskette


{ \tilde{ {\mathfrak a} }
}


{ = }{ {\mathfrak a}S
}


{  }{
}


{    }{
}


{   }{
}


}
{}{}{}
das
\definitionsverweis {Bildideal}{}{.}
Zeige, dass


\mavergleichskette


{\vergleichskette


{ {\mathfrak a}^n S
}


{ = }{ \tilde{ {\mathfrak a} }^n
}


{  }{
}


{    }{
}


{   }{
}


}
{}{}{}
ist.


}


{} {}


\inputaufgabe


{}


{


Es sei $K$ ein
\definitionsverweis {Körper}{}{.}
Wir betrachten in


\mathl{K[X,Y]}{} die beiden
\definitionsverweis {Primideale}{}{}


\mavergleichskettedisp


{\vergleichskette


{ {\mathfrak p}
}


{ =} { (X)
}


{ \subset} { (X,Y)
}


{ =} { {\mathfrak m}
}


{ } {
}


}
{}{}{.}
Zeige, dass es kein Ideal ${\mathfrak a}$ mit


\mavergleichskettedisp


{\vergleichskette


{ {\mathfrak p}
}


{ =} { {\mathfrak a} \cdot {\mathfrak m}
}


{ } {
}


{ } {
}


{ } {
}


}
{}{}{}
gibt.


}


{} {}


  \zwischenueberschrift{Aufgaben zum Abgeben}


\inputaufgabe


{3}


{


Bestimme alle Lösungen der Gleichung


\mavergleichskettedisp


{\vergleichskette


{ x^2+y^2+xy
}


{ =} { 1
}


{ } {
}


{ } {
}


{ } {
}


}
{}{}{}
für die
\definitionsverweis {Körper}{}{}


\mavergleichskette


{\vergleichskette


{ K
}


{ = }{ \mathbb F_2
}


{  }{
}


{    }{
}


{   }{
}


}
{}{}{,}
$\mathbb F_4$ und $\mathbb F_8$.

Man kann für die Körper die Darstellungen auf [[Endliche Körper/Nicht Primkörper/Einige Operationstafeln]] verwenden.


}


{} {}


\inputaufgabe


{3}


{


Es sei


\mavergleichskette


{\vergleichskette


{ M
}


{ = }{ \{ P_1,P_2 , \ldots , P_m \}
}


{ \subseteq }{ { {\mathbb A}_{ \R }^{  n   } }
}


{    }{
}


{   }{
}


}
{}{}{}
eine endliche Menge von Punkten. Zeige, dass $M$ die
\definitionsverweis {Nullstellenmenge}{}{}
eines einzigen Polynoms ist.


}


{} {}


\inputaufgabe


{3}


{


Zeige, dass die Menge der reellen
\definitionsverweis {trigonalisierbaren}{}{}
$(2\times 2)$-\definitionsverweis {Matrizen}{}{}
im


\mathl{{ {\mathbb A}_{ \R }^{  4   } }}{} keine
\definitionsverweis {affin-algebraische Menge}{}{}
ist.


}


{} {}


\inputaufgabe


{4}


{


Bestimme die Schnittpunkte der beiden Ellipsen


\mathdisp {\left\{ (x,y) \in\R^2  \mid   4x^2-3xy+2y^2 = 7         \right\} \text{   und } \left\{ (x,y) \in\R^2  \mid   3x^2+4xy+5y^2 = 8         \right\}} { . }


}


{} {}


\inputaufgabe


{4}


{


Es sei $S$ das Nullstellengebilde in ${ {\mathbb A}_{ K }^{  3   } }$, das durch die Gleichung


\mavergleichskettedisp


{\vergleichskette


{ x^2+y^2
}


{ =} { z^2
}


{ } {
}


{ } {
}


{ } {
}


}
{}{}{}
gegeben ist. Der Schnitt von $S$ mit einer Ebene $E$ ist eine Kurve und wird in $E$ durch eine Gleichung in zwei
\zusatzklammer {geeigneten} {} {}
Variablen beschrieben. Finde eine solche Gleichung für die Ebenen


\mathdisp {E_1=V( x),\, E_2=V(z-1), \, E_3 =V(x+2y+3z ),\, E_4 = V(3x-2z)} { . }


}


{} {}
